\section{Mechanization}
\label{sec:mechanization}

Three tools carry the development, at assurance levels \cref{tab:assurance}
labels claim by claim. TLC checks the witnesses, the lattice and ledger
matrices and the fifteen \cref{thm:corc} verdicts (bounded, exhaustive
within bounds). TLAPS proves \cref{thm:corc} (\TlapsCommitOrder{} obligations), a
coherence check of the chain definitions (\texttt{Hierarchy.tla},
\TlapsHierarchy{}) and a same-agent exclusion lemma
(\texttt{A1LowerBound.tla}, \TlapsLowerBound{}); none is a substantive result
except the first. Verus~\cite{Lattuada2023Verus} carries the substantive
mechanization: \VerusCurated{} obligations (\VerusFull{} counting four helper
files), zero errors, no \texttt{assume} or \texttt{admit} in any proof body,
regenerated by \texttt{verus\_count.sh} against a pinned build
(\texttt{verus-version.txt}; a newer verifier makes proof files stop
compiling rather than fail, which the counter reports as such).
Of these, \VerusLoadBearing{} bear a deployed or measured artifact's
safety claim and \VerusScaffold{} are scaffolding (superseded refinements,
the $A_4$ model, the model-only Verus mirror of \texttt{Guarded.tla}),
verified but not load-bearing; supplement \S\,S inventories every file.



\begin{table*}[t]
\centering
\caption{Every claim at its assurance level. \emph{Model}: checked by TLC
within bounds; \emph{proof}: TLAPS or Verus, no \texttt{assume}/\texttt{admit};
\emph{bytes}: the verified code is the code measured; \emph{port}: unverified
code gated against a verified artifact on the committed data; \emph{open}:
stated, not established. Correspondences name the hypothesis under which two
formalizations agree and the gate (\texttt{harness/macaudit/}) that checks it.}
\label{tab:assurance}
\scriptsize
\setlength{\tabcolsep}{4pt}
\begin{tabular}{@{}p{3.1cm}p{2.0cm}p{2.5cm}p{3.6cm}p{5.0cm}@{}}
\toprule
Claim & Definition & Proof / check & Executed by & Correspondence, hypothesis, gate; status \\
\midrule
Anomaly witnesses, lattice realizability & \TLA{} & TLC, $|A|{=}2$, $\le2$ cells, $\le2$ values, $\mathrm{MaxOps}\le4$ (witnesses), $\le3$ (matrix) & --- & \textbf{model}: bounded; \LatticeRuns{}$+$\LatticeMaxOpsThreeRuns{}$+32$ two-value safety runs \\
\cref{thm:corc} & \TLA{} & TLAPS \TlapsCommitOrder{} & --- & \textbf{proof}; hypotheses H1--H3 are model invariants (TLC) and H1 is checked on all \HypRecords{} committed records (\texttt{hypothesis\_gates}: \HypSameAgentOverlaps{} violations) \\
\texttt{detect\_a1} sound \& complete & Verus spec & Verus & scorer (port) & \textbf{proof} against the restatement; restatement $=$ \cref{def:a1} under \cref{lem:same-agent} (TLAPS step); port gated on every cell (\texttt{predicate\_matrix\_a1}); \textbf{port} not yet wired \\
\texttt{detect\_a2} & Verus spec (id, signature) & Verus & $L_4$ dispatcher, scorer & \textbf{proof}; \cref{def:a2} (identifiers) coincides under the re-sign-as-replace encoding; \textbf{encoding stated, not proved} \\
\texttt{detect\_a3} & Verus spec (unsupported read) & Verus & scorer & \textbf{proof} for the screen; $=$ \cref{def:a3} under value-unique writes, which hold on every scored trace (\texttt{hypothesis\_gates}: \HypLiveLtwoSessions{} sessions, \HypLiveLtwoCollisions{} collisions) \\
\texttt{detect\_a6} & Verus spec & Verus & scorer & \textbf{proof}; single-record equivalence to \cref{def:a6} \\
$L_1$ stores window-free & Verus & Verus (lock predicate) & pilot (bytes); live study \emph{validated} arm (application layer) and \emph{verified} arm (bytes via bridge) & \textbf{bytes} in the pilot and the verified arm; \textbf{port} in the validated arm; refinement to the abstract machine \textbf{open} (bounded simulation only) \\
$L_2$ output commit & Verus & Verus & \texttt{l2\_exec} (bytes) & \textbf{bytes}; sequential object, caller serializes \\
$L_3$ sequencing & Verus & Verus & twin (inspection) & \textbf{port} \\
$L_4$ validation at dispatch & Verus & Verus & \texttt{l4\_exec} (bytes) & \textbf{bytes} \\
\Rp{} on the live records & \cref{def:residue} & --- & \texttt{a7\_rescore} (Python) & \textbf{port}, unverified; independent of the outcome column; extractor recall \HypBasisRecovered/\HypBookings{} \\
Basis ledger excludes \Rp{} & \TLA{} $G_7$; Verus model & TLC \LatticeThirtyTwoRuns{} runs; Verus \LfiveVerusIntro & probe 303 ledger arm (Python) & \textbf{model} and \textbf{proof}; executable ledger is a \textbf{port} (harness code) \\
\bottomrule
\end{tabular}
\end{table*}

\subsection{Verified detectors, and one transcription gap}
\label{sec:verus-detectors}
The four detectors are verified sound \emph{and} complete against Verus
\texttt{spec} restatements of the \TLA{} predicates (\VerusDetector{}
obligations). \texttt{detect\_a1} is a two-sided equivalence whose
completeness half strengthens each of three nested loop invariants to carry
``no witness in the prefix searched.'' \texttt{detect\_a2} and
\texttt{detect\_a6} are single-record equivalences; \texttt{detect\_a3} is
exact for the flat-trace residue of \cref{sec:a3}.

The restatement is hand-written and checked by inspection, and inspection
found a gap we state rather than hide. The Verus \texttt{OpRecord} carries no
agent field, so its \A1 predicate omits the different-agent conjunct of
\cref{def:a1}; and it decides ``no witness'' when a read value is absent,
where the \TLA{} compares against $\mathrm{NULL}$. Neither matters on the
histories the model admits, and we say why.

\begin{lemma}[Same-agent exclusion]
\label{lem:same-agent}
Under per-agent non-overlap and read-before-commit, if
$c\in h[i].\rs\cap h[k].\ws$ and $h[i].\rt<h[k].\wt<h[i].\wt$, then
$h[i].\fld{agent}\neq h[k].\fld{agent}$.
\end{lemma}
The lemma is step $\langle2\rangle2$ of the TLAPS proof of \cref{thm:corc}
and the content of \texttt{A1LowerBound.tla}. It makes the Verus predicate
and \cref{def:a1} coincide on every history satisfying the model's invariants;
on a history that violates non-overlap the Verus predicate is strictly
stronger---it over-reports, the safe direction for a detector. The
instrumentation records a value for every read cell, so the absent-value case
does not arise: across all \HypRecords{} committed operation records there
is no read cell without a value and \HypSameAgentOverlaps{} same-agent
overlapping windows (\texttt{hypothesis\_gates.py --check}). The published tables were
scored by \texttt{rust-analyser}, a hand-written port of the verified
detector; a gate (\texttt{predicate\_matrix\_a1.py}) scores every committed
cell under the \TLA{}, the Verus and the port's transcription and fails on
any divergence, and all three agree on every cell. Wiring the verified
\texttt{detect\_a1} into the scorer is a build step the artifact's delegation
shim prepares and this version has not taken.

\subsection{$L_1$: two concurrent stores whose lock is verified, not assumed}
\label{sec:l1}
Three abstract state machines---pessimistic locking, serializable snapshot
isolation (SSI), and default snapshot isolation with its no-write bypass---are
proved \A1-free under sequential interleaving ($40$ obligations; default-SI
conditionally, on workloads with no read-only commit, which is the residual
\cref{sec:pilot} measures). The deployed stores do not rest on these: each is
an exec-mode store whose only access path is \texttt{vstd::rwlock::RwLock},
the reader--writer lock implemented and verified inside the Verus standard
library, carrying the store's invariant as the lock predicate
(\texttt{lib\_si\_concurrent.rs}, $18$; \texttt{lib\_pess\_concurrent.rs},
$31$; no axiom, no \texttt{external\_body}). The invariant entails the absence
of any stale-generation window in the store's own record stream, under any
interleaving of threads. Theorem L1j (\texttt{lib\_si\_projection.rs}, $3$)
transports that guarantee to the detector's history for \emph{faithful}
emitters---one \texttt{OpRecord} per store record, cells a subset of the
record's. The default-SI \texttt{tick} path is the standing counterexample to
faithfulness and is exactly the $\PilotSiEditResidual\%$/$\PilotSiTriageResidual\%$
residual of \cref{tab:pilot}. Four earlier refinements ($86$ obligations,
two axioms) relate the abstract machines to the representations the runtime
shipped \emph{before} the verified-lock rebuild and bear no deployed claim;
refining the two concurrent stores to the abstract machines is the open
proof step (a bounded simulation over $506{,}880$ transitions finds no
counterexample; supplement \S\,H).

\subsection{$L_2$: output commit}
\label{sec:l2}
The $L_2$ model (\texttt{lib\_l2\_safety.rs}, $33$) carries per transaction a
predecessor set recorded as a causal closure at read time. A commit is valid
iff the read set is fresh against the committed store and every predecessor
is committed and unaborted; an abort cascades to every transaction whose
closure holds the aborter; effects leave by a separate \texttt{externalize}
step enabled only once the closure has externalized; an externalized
transaction cannot abort. Theorem $L_2$h: no state satisfying the inductive
invariant $\fld{inv\_output\_commit}$ contains a witness of \cref{def:a3}
(\texttt{lemma\_l2\_reachable\_no\_a3}), under arbitrary interleaving of the
six steps; the closure invariant makes the one-level cascade provably
sufficient for the transitive case. Theorem $L_2$i: every committed read has
a surviving producer of exactly its value, and (temporal strengthening) that
producer committed no later than the read. Theorem $L_2$j
(\texttt{lib\_l2\_projection.rs}, $3$): a faithful emission of an
$L_2$-supported state contains no residue witness, by a bridge that never
refers to a read's producer and so cannot assume what it proves. The exec
runtime (\texttt{lib\_l2\_exec.rs}, $73$ end to end, $40$ net of the
re-included model) refines the model per transition with zero added axioms
(its \texttt{u64} interning is natively injective); every mutating entry
point checks its own guard and reports whether it acted, so the guarantee
survives proof erasure under plain \texttt{cargo}, and the file compiled by
the reference runtime is byte-identical.

\subsection{$L_3$: sequencing, and $L_4$: validation at dispatch}
\label{sec:l3}\label{sec:l4}
Theorem $L_3$g (\texttt{lib\_l3\_sequencer.rs}, $5$): a commit-order
sequencer that externalizes effect $k$ once effects $0..k$ have completed
yields an externalization time monotone in issuance index for \emph{every}
completion schedule, hence no \A6, without forbidding concurrency; a
reversed schedule exhibits \A6 unsequenced and none sequenced. The saga
model (\texttt{lib\_l3\_safety.rs}, $6$) is the strict-serialization special
case and composes with $L_2$ on disjoint carriers. The exec sequencer
(\texttt{lib\_l3\_exec.rs}, $7$) is verified; its measurement
(\cref{sec:a6-eval}) runs a dependency-free twin, a gap the trust base
records.

Resolving a planned tool against a registry snapshot does not prevent \A2: a
snapshot changes what an operation resolves, not the tool its call reaches
(Theorem $L_4$d, a concrete execution). The $L_4$ model
(\texttt{lib\_l4\_safety.rs}, $9$) reads \A2 from a \emph{dispatch record}
and validates at dispatch---place the call only if the planned tool is still
live with the pinned signature; registry churn interleaves arbitrarily and
cannot create \A2 because \A2 is a property of the dispatch. The exec
dispatcher (\texttt{lib\_l4\_exec.rs}, $11$) refuses a revoked or re-signed
tool and dispatches an untouched one, and \cref{sec:l4-eval} measures those
bytes. Validation and dispatch are one atomic step of a registry the runtime
owns; a remote tool server (MCP~\cite{AnthropicMCP2024}) offers no such
atomicity and the check-then-call gap there is narrowed, not closed.


\subsection{The basis ledger}
\label{sec:ledger}
The ledger $\ell$ answers \Rp{} and composes with any discipline; the live
study runs it over $L_1$ (\cref{tab:live}). Its model
(\texttt{lib\_l5\_ledger.rs}, \LfiveVerus) carries the
store's committed values and, per effect, its basis map, an externalized
flag and an obligation flag. Three steps: \emph{register} commits an effect
with its basis (released at once, or held under output commit; if released
at once on a basis the store already disagrees with, the obligation is set
at commit); \emph{write} installs a value and marks, in the same step, every
externalized effect whose basis holds that cell with a different value;
\emph{externalize} releases a held effect and marks it if its basis is stale
by then. The invariant $\mathit{inv\_l5}$---every externalized effect with a
stale basis carries an obligation---is preserved by each step (the write
step is the load-bearing case: the cell that became stale is the cell the
step marked, and any other stale cell was stale before), and the ledger theorem:
no state satisfying $\mathit{inv\_l5}$ contains a witness of \cref{def:residue}.
The model's witness is the state form---an externalized, unowed effect whose
basis the store now disagrees with---which contains \cref{def:residue}'s (a
latest post-commit writer that differs implies a current value that
differs), so excluding it excludes the definition's witness and, in
addition, an effect released on a basis already stale at commit.
Non-vacuity in both directions: the unguarded write reaches \Rp{} in two
steps from an effect derived from an unwritten cell, and under the ledger the
same write registers the obligation while a write that restores the value
leaves no witness. The ledger refuses nothing; refusing to release a
stale-basis effect, or holding it for a deadline, is a policy the executable
runtime layers on top, and what a registered obligation achieves in the
world is measured in \cref{sec:live-a1}. The model is \LfiveVerusIntro{}
under the pinned build; its obligations are counted as scaffolding
(\VerusScaffold{} with it) because the executable ledger of the live study is
harness code, not these bytes. An exec ledger on the $L_2$/$L_4$ pattern is
not in this version.

\subsection{Composition}
\label{sec:composition}
\texttt{lib\_consistency\_lattice.rs} takes the product of the $L_2$, $L_3$
and $L_4$ states over their disjoint carriers and proves that the conjunction
of the three invariants excludes \A1 (image), \A3, \A6 and \A2 jointly, with
three non-interference lemmas; $53$ obligations, of which $48$ are the
re-verified component modules. The ledger's carrier (basis, obligation) is
disjoint from all three, and adding it to the product is the same
construction; it is not in this version.

\subsection{Trust base}
\label{sec:trust}
``Zero \texttt{assume}, zero \texttt{admit}'' is a statement about proof
bodies. What is trusted, in full: (i) Verus, its SMT backend and
\texttt{vstd}, including \texttt{group\_hash\_axioms} for the
\texttt{HashMap}-backed exec states; (ii) two axioms in the four superseded
refinements (string injectivity, null sentinel) and (iii) three
\texttt{external\_body} stubs for a \texttt{parking\_lot::Mutex}, neither
bearing a deployed claim since the verified-lock rebuild; (iv) the hand
restatement of the \TLA{} predicates in Verus, with the gap of
\cref{sec:verus-detectors} stated and gated; (v) the hand-written scorer that
produced the tables, gated against the verified predicate on every committed
cell; (vi) for $L_3$, the inspection-level correspondence between the exec
artifact and the measured twin; (vii) the \texttt{pub} fields of the exec
runtimes, which a caller can write around the API; (viii) the marshalling
that flattens the version map into the verified validation gate, and the
interning bijection and abort counter outside the verified store; (ix) the
Python and Rust instrumentation that produces every \texttt{OpRecord},
including the read sets it declares---a detector is as good as the read set
it is handed, and recovering implicit read sets is a separate
problem~\cite{Khan2026SBus}; (x) the provenance correspondence between the
two formalizations of \A3, which agree only under value-unique writes
(checked on every scored trace); (xi) in the live study's verified arm, the
$100$-line JSON bridge and the harness's interning of cells and values. No
safety-bearing proof invokes a writes-as-function-of-reads assumption, and
the development carries no probabilistic obligation.
